\addchap{MỞ ĐẦU}\label{chap:introduction}

\addsec{Tổng quan tình hình nghiên cứu trong nước và ngoài nước}

Trong các tình huống thiên tai bão lũ, thông tin cầu cứu thường được gửi từ nhiều nguồn, tại những thời điểm khác nhau và dưới nhiều hình thức như văn bản, hình ảnh và vị trí địa lý. Mỗi loại dữ liệu phản ánh một khía cạnh của tình huống: hình ảnh cung cấp thông tin về mức độ ngập lụt và tình trạng hiện trường, văn bản mô tả nhu cầu hỗ trợ, trong khi vị trí địa lý giúp xác định khu vực cần tiếp cận. Tuy nhiên, các báo cáo có thể bị trùng lặp do được gửi nhiều lần hoặc do nhiều người cùng phản ánh một sự kiện. Ngược lại, những tình huống cứu hộ khác nhau có thể xảy ra tại các vị trí gần nhau. Vì vậy, một hệ thống hỗ trợ cứu hộ cần có khả năng phân tích từng báo cáo, nhận diện các thông tin liên quan và tổng hợp chúng thành những sự kiện phục vụ quá trình tiếp nhận, xác minh và điều phối.

Tại Việt Nam, các nghiên cứu ứng dụng trí tuệ nhân tạo trong phòng, chống thiên tai đã được triển khai theo nhiều hướng, tiêu biểu như ứng dụng học máy trong dự báo lũ, ước lượng mức nước ngập từ hình ảnh và xây dựng bản đồ phân vùng nguy cơ lũ lụt. Những nghiên cứu này cho thấy tiềm năng của học máy, học sâu và xử lý dữ liệu không gian trong việc hỗ trợ giám sát, đánh giá tình hình thiên tai. Tuy nhiên, phần lớn các hướng tiếp cận tập trung vào dự báo, nhận diện hoặc đánh giá hiện trạng ngập lụt. Bài toán tiếp nhận, phân tích và tổng hợp các yêu cầu cứu hộ trong điều kiện hạ tầng mạng không ổn định vẫn đặt ra nhiều vấn đề cần nghiên cứu. Khác với các hệ thống dự báo khí tượng hoặc thủy văn, đề tài tập trung vào quá trình xử lý thông tin sau khi yêu cầu cứu hộ được người dùng tạo ra, nhằm hỗ trợ việc xác định và quản lý các tình huống cần xử lý.

Trên thế giới, việc khai thác dữ liệu từ mạng xã hội và các nguồn thông tin cộng đồng đã mở ra nhiều hướng nghiên cứu phục vụ ứng phó thiên tai. Sakaki và cộng sự nghiên cứu khả năng sử dụng thông tin do người dùng tạo ra để phát hiện sự kiện theo thời gian \citep{sakaki2010earthquake}. Bộ dữ liệu CrisisMMD cung cấp nguồn dữ liệu đa phương thức phục vụ nghiên cứu phân tích thông tin thiên tai \citep{alam2018crisismmd}. Bên cạnh đó, FloodNet hỗ trợ nghiên cứu nhận diện và phân tích cảnh ngập lụt từ ảnh chụp trên không \citep{rahnemoonfar2021floodnet}. Các công trình này cung cấp cơ sở tham khảo cho việc xây dựng những phương pháp phân tích dữ liệu hình ảnh và văn bản trong bối cảnh thiên tai. Đối với mô hình phân loại ảnh của đề tài, bộ dữ liệu được sử dụng trực tiếp trong huấn luyện là \texttt{Dataset\_Flood} do nhóm xây dựng. FloodNet và CrisisMMD được sử dụng làm tài liệu tham khảo; chưa có bằng chứng cho thấy hai bộ dữ liệu này đã được tích hợp trực tiếp vào tập huấn luyện của mô hình hiện tại.

Song song với nghiên cứu phân tích dữ liệu đa phương thức, trí tuệ nhân tạo tại biên (Edge AI) ngày càng được quan tâm nhờ khả năng thực hiện suy luận trực tiếp trên thiết bị đầu cuối. Cách tiếp cận này giúp giảm sự phụ thuộc vào máy chủ và hạn chế lượng dữ liệu cần truyền qua mạng, đặc biệt phù hợp với những môi trường có kết nối không ổn định. MobileNetV3 là một kiến trúc mạng nơ-ron được thiết kế nhằm tối ưu hiệu quả tính toán trên thiết bị di động \citep{howard2019mobilenetv3}. Trong nghiên cứu này, kiến trúc MobileNetV3-Large được sử dụng để xây dựng mô hình phân loại hình ảnh thành bốn mức biểu diễn tình trạng ngập lụt. Các kỹ thuật chuyển đổi định dạng, lượng tử hóa và cắt tỉa mô hình được khảo sát nhằm đánh giá khả năng giảm chi phí tính toán và hỗ trợ triển khai trên thiết bị có tài nguyên hạn chế. Kết quả thực nghiệm trên CPU máy tính được phân tích riêng với kết quả trên thiết bị Android do sự khác biệt về phần cứng, môi trường thực thi và chi phí tiền xử lý dữ liệu.

Đối với bài toán tổng hợp các yêu cầu cứu hộ, thuật toán phân cụm và phát hiện cộng đồng có thể được sử dụng để xác định những báo cáo có liên quan đến cùng một sự kiện. Louvain và Leiden khai thác cấu trúc cộng đồng trong đồ thị để tìm các nhóm đối tượng có mức độ liên kết cao \citep{blondel2008louvain,traag2019leiden}. Trong khi đó, DBSCAN là phương pháp phân cụm dựa trên mật độ, có thể sử dụng làm phương pháp đối chứng trong thực nghiệm \citep{ester1996dbscan}. Tuy nhiên, chất lượng phân cụm không chỉ phụ thuộc vào thuật toán mà còn chịu ảnh hưởng của cách xác định độ tương đồng, thông tin không gian, thời gian và khả năng xử lý các báo cáo thiếu dữ liệu. Một phương pháp có chỉ số phân cụm tốt vẫn có thể tạo ra những nhóm sự kiện không phù hợp với yêu cầu điều phối, chẳng hạn như gộp các tình huống độc lập hoặc tạo ra điểm đến không tương ứng với sự kiện thực tế. Do đó, việc đánh giá cần xem xét cả chất lượng phân hoạch và ảnh hưởng của sai lệch phân cụm đến quá trình điều phối.

Trong công tác cứu trợ thiên tai, việc phân bổ nguồn lực cứu hộ cần cân nhắc đồng thời hiệu quả, tính công bằng và khả năng đáp ứng các nhu cầu khẩn cấp \citep{gralla2014review}. Từ những vấn đề trên, nghiên cứu không chỉ tập trung vào phân loại và phân cụm thông tin mà còn xem xét quá trình xếp hạng mức độ ưu tiên và hệ quả của kết quả phân cụm đối với việc lập lịch điều phối. Ba câu hỏi nghiên cứu được đặt ra nhằm so sánh các phương pháp phân cụm đồ thị (RQ1), đánh giá tính ổn định của điểm ưu tiên trước dữ liệu trùng lặp (RQ2) và khảo sát ảnh hưởng của sai lệch phân cụm đến kết quả điều phối (RQ3).

Các kết quả thực nghiệm cũng cho thấy những giới hạn của các phương pháp được khảo sát. Phương pháp phân cụm sử dụng cổng không gian dạng tích chưa chứng minh được ưu thế tổng quát so với phương pháp dạng cộng. Điểm ưu tiên có chặn chưa cho thấy sự phù hợp nhất quán với lợi ích được xác định độc lập trong mô phỏng. Đồng thời, những đặc tính toán học của từng thành phần chưa đủ để khẳng định khả năng cải thiện kết quả điều phối. Các kết quả này giúp làm rõ điều kiện áp dụng của từng phương pháp và những vấn đề cần tiếp tục nghiên cứu khi phát triển hệ thống hỗ trợ cứu hộ.


\addsec{Lý do chọn đề tài}

Việt Nam là quốc gia thường xuyên chịu ảnh hưởng của bão lũ và các hiện tượng thời tiết cực đoan. Trong những tình huống thiên tai, việc tiếp nhận và xử lý thông tin cầu cứu kịp thời có ý nghĩa quan trọng đối với hoạt động tìm kiếm, cứu hộ và cứu nạn. Tuy nhiên, thông tin thường được gửi qua nhiều kênh khác nhau, có thể thiếu vị trí chính xác, không đầy đủ về tình trạng hiện trường hoặc bị trùng lặp. Khi số lượng báo cáo tăng, việc tổng hợp thủ công gây khó khăn cho quá trình xác định mức độ khẩn cấp và nhận diện các khu vực thực sự cần hỗ trợ.

Trong quá trình tiếp nhận thông tin, người điều phối phải đọc và kiểm tra từng báo cáo, xác định vị trí, đánh giá nhu cầu cứu hộ và xem xét liệu thông tin mới có liên quan đến sự kiện đã được ghi nhận hay không. Các báo cáo trùng lặp có thể khiến mức độ tập trung của yêu cầu cứu hộ tại một khu vực bị đánh giá cao hơn thực tế. Ngược lại, những báo cáo thiếu thông tin vị trí có nguy cơ bị bỏ sót nếu hệ thống chỉ tập trung hiển thị các điểm có tọa độ trên bản đồ. Vì vậy, việc xây dựng một quy trình tiếp nhận có khả năng lưu trữ, tổng hợp và theo dõi trạng thái báo cáo là cần thiết để hỗ trợ quá trình xác minh thông tin.

Một khó khăn khác là sự phụ thuộc vào kết nối Internet của các hệ thống xử lý tập trung. Trong điều kiện bão lũ, hạ tầng viễn thông có thể bị quá tải hoặc gián đoạn, khiến việc truyền tải hình ảnh và các dữ liệu dung lượng lớn gặp nhiều hạn chế. Nếu ứng dụng chỉ ghi nhận báo cáo sau khi gửi thành công, thông tin cầu cứu có thể bị mất khi thiết bị mất kết nối. Điều này đặt ra yêu cầu cần có cơ chế lưu trữ cục bộ, duy trì dữ liệu trong thời gian ngoại tuyến và đồng bộ khi kết nối được khôi phục.

Để giải quyết vấn đề trên, đề tài lựa chọn hướng tiếp cận kết hợp Edge AI với cơ chế lưu trữ và truyền dữ liệu tối giản. Thông tin được ghi nhận và lưu tại thiết bị trước khi gửi, trong khi mô hình AI hỗ trợ phân tích hình ảnh để tạo ra các thông tin mô tả ngắn gọn. Tùy theo điều kiện kết nối, hệ thống có thể gửi dữ liệu tối thiểu hoặc sử dụng ảnh đã nén nhằm giảm lượng dữ liệu truyền tải. Khi kết nối được khôi phục, những báo cáo chưa được gửi thành công sẽ tiếp tục được đồng bộ về máy chủ. Kết quả suy luận tại thiết bị đóng vai trò hỗ trợ phân tích, không thay thế hoàn toàn thông tin do người dùng cung cấp hoặc quá trình xác minh tại trung tâm.

Bên cạnh khả năng hoạt động khi mạng không ổn định, việc tổng hợp các báo cáo trùng lặp và xác định thứ tự ưu tiên cứu hộ cũng là vấn đề cần nghiên cứu. Những tiêu chí như mức độ ngập, số người cần hỗ trợ và sự hiện diện của các đối tượng dễ bị tổn thương có thể được sử dụng để xây dựng điểm ưu tiên. Tuy nhiên, một công thức có cấu trúc hợp lý chưa đủ để chứng minh thứ tự ưu tiên phản ánh chính xác nhu cầu cứu hộ. Tương tự, việc hạn chế ảnh hưởng của bản sao ở tầng tính điểm không đồng nghĩa với khả năng chống chịu trước dữ liệu trùng lặp ở tầng phân cụm. Do đó, cần đánh giá các phương pháp trên bằng những thực nghiệm có đối chứng thay vì chỉ dựa vào tính chất lý thuyết.

Xuất phát từ những vấn đề thực tiễn và kỹ thuật nêu trên, đề tài ``Xây dựng hệ thống phân tích đa phương thức và phân cụm sự kiện cứu hộ bão lũ dựa trên Edge AI'' được thực hiện nhằm nghiên cứu khả năng kết hợp phân tích dữ liệu tại thiết bị, truyền thông tin trong điều kiện mạng không ổn định và phân cụm các báo cáo cứu hộ. Hệ thống hướng đến việc hỗ trợ tiếp nhận, tổng hợp và trực quan hóa thông tin, đồng thời cung cấp cơ sở để nghiên cứu các phương pháp xếp hạng ưu tiên và hỗ trợ điều phối nguồn lực trong những tình huống thiên tai.


\addsec{Mục tiêu đề tài}

Mục tiêu chung là xây dựng hệ thống ứng dụng trí tuệ nhân tạo trên thiết bị di động để tiếp nhận, phân tích thông tin và hỗ trợ ra quyết định cứu hộ bão lũ trong điều kiện mạng không ổn định, hướng tới giảm thiểu thiệt hại về người và tài sản.

Các mục tiêu cụ thể bao gồm:

\begin{enumerate}
    \item Nâng cao khả năng nhận diện và phân tích tình huống cứu hộ từ văn bản và hình ảnh.
    \item Duy trì hoạt động khi mạng yếu hoặc gián đoạn bằng xử lý tại thiết bị và truyền dữ liệu tối giản.
    \item Đánh giá mức độ khẩn cấp và gom nhóm các tình huống cứu hộ để hỗ trợ ra quyết định, điều phối nguồn lực.
\end{enumerate}

Đối với phần nghiên cứu thuật toán, ba câu hỏi nghiên cứu được đặt ra như sau:

\begin{enumerate}
    \item RQ1: Phương pháp phân cụm đồ thị sử dụng cổng không gian dạng tích có cải thiện chất lượng phân cụm so với phương pháp dạng cộng khi các cấu hình được lựa chọn độc lập hay không?
    \item RQ2: Phương pháp tính điểm ưu tiên tổng hợp theo họ bằng chứng có duy trì tính ổn định trước các báo cáo trùng lặp và phản ánh phù hợp lợi ích được xác định độc lập trong mô phỏng hay không?
    \item RQ3: Những sai lệch trong quá trình phân cụm, bao gồm tách sai, gộp sai và hình thành điểm đến giả, ảnh hưởng như thế nào đến kết quả lập lịch điều phối cứu hộ?
\end{enumerate}

Các mục tiêu được đánh giá bằng những chỉ số phù hợp với từng bài toán. Đối với mô hình phân loại ảnh, Accuracy và macro-F1 được sử dụng để đo chất lượng dự đoán. Chỉ số ARI đánh giá mức độ phù hợp của phân hoạch, trong khi NDCG@5 được sử dụng để đánh giá chất lượng xếp hạng ưu tiên. Đối với mô hình điều phối mô phỏng, các chỉ số harm và tỷ lệ trễ hạn phản ánh những khía cạnh khác nhau của kết quả lập lịch. Riêng phần mềm được kiểm tra thông qua các ca thử nghiệm về tạo, lưu trữ, gửi và đồng bộ báo cáo.

Các kết quả đánh giá được phân tích theo từng tầng chức năng. Chất lượng phân loại ảnh, phân cụm sự kiện hoặc xếp hạng ưu tiên không được xem là bằng chứng trực tiếp cho hiệu quả cứu hộ thực tế nếu chưa có những thử nghiệm tương ứng.


\addsec{Phương pháp nghiên cứu}

Đề tài kết hợp nghiên cứu tài liệu, xử lý dữ liệu, thiết kế hệ thống và so sánh thực nghiệm qua các bước sau:

\begin{enumerate}
    \item \textbf{Thu thập và xử lý dữ liệu:} Khảo sát các nghiên cứu về nhận diện ảnh, phân tích văn bản, Edge AI và phân cụm; chuẩn hóa, gán nhãn dữ liệu phục vụ huấn luyện và đánh giá.
    \item \textbf{Phân tích yêu cầu và thiết kế hệ thống:} Xác định chức năng, yêu cầu hiệu năng và khả năng hoạt động khi mạng gián đoạn; thiết kế kiến trúc, cơ sở dữ liệu và luồng trao đổi giữa ứng dụng với máy chủ.
    \item \textbf{Huấn luyện và tối ưu mô hình:} Huấn luyện MobileNetV3-Large; so sánh chuyển đổi định dạng, lượng tử hóa và cắt tỉa trên cùng tập kiểm thử theo độ chính xác, kích thước và độ trễ. Kết quả CPU máy tính và điện thoại được phân tích riêng.
    \item \textbf{Phân cụm và xếp hạng:} Xây dựng đồ thị từ vị trí, thời gian và ngữ cảnh; so sánh chất lượng phân cụm, độ ổn định điểm ưu tiên và kết quả điều phối mô phỏng khi có dữ liệu thiếu, nhiễu hoặc trùng lặp. Cấu hình được chọn trên tập hiệu chỉnh trước khi đánh giá trên tập kiểm thử; các phương pháp dùng cùng điều kiện nhiễu trong mỗi lần chạy.
    \item \textbf{Phát triển và kiểm thử phần mềm:} Xây dựng ứng dụng Flutter, backend Python và dashboard; kiểm tra tạo, lưu, gửi, đồng bộ báo cáo và khả năng truyền dữ liệu dưới mạng yếu giả lập.
    \item \textbf{Đánh giá kết quả:} Đối chiếu dự đoán, cấu hình và bảng số liệu đã lưu; sử dụng bootstrap, kiểm định Wilcoxon và hiệu chỉnh Holm cho các so sánh ghép cặp theo seed. Đơn vị đánh giá là ảnh, lần chạy RQ1, seed RQ2/RQ3 hoặc ca kiểm thử tương ứng với từng bài toán.
\end{enumerate}

Thông tin đầu vào thuật toán trong \texttt{algorithm\_input.json} được tách khỏi nhãn sự kiện, lợi ích mô phỏng, thời hạn xử lý và mức thiệt hại dùng để đánh giá. Riêng RQ2 sử dụng nhóm sự kiện tham chiếu để đánh giá riêng tầng xếp hạng, không coi đây là thông tin sẵn có khi vận hành. Việc tổng hợp, đối chiếu các kết quả đã lưu không được tính là một lần huấn luyện hoặc thử nghiệm thiết bị mới.

\addsec{Đối tượng nghiên cứu}

\begin{itemize}
    \item Dữ liệu ảnh ngập lụt bốn lớp, thông tin cứu hộ tiếng Việt và các trường vị trí, thời gian, mức độ khẩn cấp; báo cáo mô phỏng phục vụ phân cụm, xếp hạng và điều phối.
    \item Mô hình MobileNetV3-Large, các định dạng PyTorch, ONNX, ExecuTorch và kỹ thuật tối ưu cho điện thoại; mô hình đánh giá khẩn cấp từ 32 trạng thái có cấu trúc, không phải tập tin nhắn cứu hộ thực tế.
    \item Phương pháp xây dựng đồ thị tương đồng, phân cụm sự kiện và tính điểm ưu tiên cứu hộ.
    \item Ứng dụng Flutter, backend, cơ sở dữ liệu SQLite, dashboard và cơ chế lưu trữ, truyền, đồng bộ thông tin.
\end{itemize}

Hệ thống hướng đến người dân gửi yêu cầu và lực lượng cứu hộ; nghiên cứu chưa khảo sát thực địa với các nhóm người dùng này.

\addsec{Phạm vi nghiên cứu}

\begin{itemize}
    \item \textbf{Thời gian:} 06 tháng, từ tháng 05/2026 đến tháng 10/2026.
    \item \textbf{Nội dung và nền tảng:} Phân tích thông tin cứu hộ bằng AI tại thiết bị; tạo, lưu và gửi báo cáo qua HTTP hoặc SMS gateway; đồng bộ khi có mạng; phân cụm, xếp hạng và hiển thị trên ứng dụng di động, web. Phân cụm xử lý các tập báo cáo hiện có, dashboard cập nhật theo chu kỳ; quyết định điều phối vẫn do con người thực hiện.
    \item \textbf{Dữ liệu:} Ảnh ngập lụt đã chuẩn hóa, thông tin mô tả nhu cầu cứu hộ và các trường liên quan. RQ1 dùng dữ liệu bán tổng hợp Generator 3.0; RQ2/RQ3 dùng dữ liệu tổng hợp Candidate 4.1. Đây là các bộ dữ liệu riêng, không phải báo cáo cứu hộ thực tế hoặc một chuỗi kiểm chứng trên cùng tập mẫu.
    \item \textbf{Phân tích ảnh và văn bản:} Đánh giá riêng mô hình nhận diện mức ngập và các phương án tối ưu trên CPU, điện thoại; không khái quát hiệu năng cho mọi thiết bị. Xử lý văn bản hiện dùng từ khóa ở backend và Logistic Regression từ các dấu hiệu có cấu trúc, chưa thay thế mô hình phân tích trực tiếp tin nhắn cứu hộ tiếng Việt.
    \item \textbf{Mạng và giới hạn ứng dụng:} Tập trung lưu báo cáo và gửi lại khi có kết nối với máy chủ. DTN qua thiết bị trung gian, phương tiện mang dữ liệu hoặc UAV mới là định hướng, chưa được kiểm chứng. Nghiên cứu không đi sâu vào dự báo khí tượng--thủy văn hoặc điều hành cứu hộ tự động.
\end{itemize}

Các thực nghiệm đánh giá từng thành phần trong điều kiện đã xác định. Hiệu quả về thời gian phản ứng, giảm thiệt hại và khả năng triển khai trong thiên tai cần được kiểm chứng thêm qua thử nghiệm thực địa với đơn vị chuyên môn.
